\chapter{Introduction}

\section{Background}
The Internet of Things (IoT) connects physical objects, sensors, and various devices through communication networks, enabling data-driven services in domains such as healthcare, industrial automation, and smart environments. Alongside the expansion of IoT deployments, the number of connected components exposed to cybersecurity threats increases correspondingly~\cite{atzori2010}. The Message Queuing Telemetry Transport (MQTT) protocol is commonly used for lightweight publish-subscribe communication in IoT environments. However, MQTT-based systems can become targets for denial-of-service (DoS) and distributed denial-of-service (DDoS) attacks. The DoS/DDoS-MQTT-IoT dataset, collected from a physical IoT testbed, provides normal MQTT traffic alongside ten DoS attack scenarios to facilitate the study of attacks in this context~\cite{alatram2023}. Intrusion Detection Systems (IDS) monitor network activity to identify traffic or behavior exhibiting malicious signs. Machine learning and deep learning methods have been utilized and researched for this task due to their ability to learn characteristic patterns directly from network traffic. Their effectiveness depends partly on the availability and representativeness of the training data~\cite{he2009}. Federated Learning (FL) allows multiple clients to collaboratively train a shared model while keeping their data local; a central server distributes the global model and aggregates updates from clients across training rounds~\cite{mcmahan2017}. This approach is highly suitable when IoT data is distributed across multiple devices or locations, although communication and computational costs remain critical factors to consider in resource-constrained environments~\cite{shi2016, imteaj2022}. Keeping data local also does not equate to completely eliminating privacy concerns arising from the exchange of model updates~\cite{zhu2019}.

Privacy-preserving synthetic data generation is a potential solution to address the shortage of training data for rare attack classes in IoT networks while still ensuring strict security guarantees. Rather than relying on deep neural networks that require substantial computational resources and face convergence difficulties on edge devices~\cite{behera2022fedsyn}, the probabilistic graphical model paradigm (Probabilistic Graphical Models) approaches the problem by decomposing tabular data into low-dimensional marginal frequency distribution slices (low-dimensional marginals)~\cite{zhang2017privbayes}. Through the Chow-Liu dependency tree structure (Chow-Liu Tree)~\cite{chow1968}, this approach effectively captures the core correlations between traffic features and attack labels without being affected by the curse of dimensionality. When combined with Differential Privacy (DP)~\cite{zhang2017privbayes} on contingency tables and consistent graphical inference mechanisms (such as Private-PGM), the system can reconstruct a valid global distribution to generate conditional traffic data, establishing a foundation to enhance the performance of downstream intrusion detection systems (IDS)~\cite{mckenna2019pgm, mckenna2024aim}.
\section{Problem Statement}
\section{Research Objectives}
\section{Significance of the Study}
\section{Scope and Limitations}
\section{Thesis Structure}
